% GENERATED FILE. Edit canonical lessons, then run npm run generate:books.
% canonical-source-hash: fnv1a64:6b291bc6d794b9bd
% canonical-lessons: SA-C211-vyayamah, SA-C211-yogah, SA-C211-sramah, SA-C211-klantih, SA-C211-snanam, SA-R211-first-pass-words-from-chapters-203-206, SA-R211-second-pass-words-from-chapters-207-211

\chapter{Exercise, Yoga, Toil, Tiredness, Bath}
\label{ch:sa-exercise-yoga-toil-tiredness-bath}
\hlchaptermodality{211}

\begin{chapteropening}
\textbf{By the end of this chapter:} \emph{I can say exercise, yoga, toil, tiredness and a bath.}

Complete the last lesson of chapter 211: I can say exercise, yoga, toil, tiredness and a bath.
\end{chapteropening}

\section[\texorpdfstring{\sk{व्यायामः} (vyāyāmaḥ)}{व्यायामः (vyāyāmaḥ)}]{\sk{व्यायामः} (vyāyāmaḥ) — exercise}

\label{lesson:SA-C211-vyayamah}



\begin{quote}
Before the new one: say the Sanskrit for a classmate, then the Sanskrit for a doctor.
\end{quote}

\subsection*{You'll want to know: \sk{व्यायामः}}
\textbf{\sk{व्यायामः}} — \emph{vyāyāmaḥ} — \textquotedblleft{}exercise\textquotedblright{}.

\begin{grammarlens}[title={What you've built}]
One more word for this chapter.
\end{grammarlens}

\subsection*{Your turn}
\begin{itemize}
\raggedright
  \item \emph{Say it:} \emph{vyāyāmaḥ}
  \item \emph{Say it:} \emph{vyāyāmaḥ}, once more
  \item \emph{Say it:} say \emph{vyāyāmaḥ}
  \item \emph{Recall:} say the Sanskrit for a classmate, then the Sanskrit for a doctor, then say \emph{vyāyāmaḥ} again
\end{itemize}

\begin{tcolorbox}[breakable,colback=teal!4,colframe=teal!35!black,title={Before you move on}]
What does \sk{व्यायामः} mean? (\textquotedblleft{}Exercise\textquotedblright{}.) Say it once more.
\end{tcolorbox}
\section[\texorpdfstring{\sk{योगः} (yogaḥ)}{योगः (yogaḥ)}]{\sk{योगः} (yogaḥ) — yoga}

\label{lesson:SA-C211-yogah}



\begin{quote}
Before the new one: say the Sanskrit for a doctor, then the Sanskrit for exercise.
\end{quote}

\subsection*{You'll want to know: \sk{योगः}}
\textbf{\sk{योगः}} — \emph{yogaḥ} — \textquotedblleft{}yoga\textquotedblright{}.

\begin{grammarlens}[title={What you've built}]
One more word for this chapter.
\end{grammarlens}

\subsection*{Your turn}
\begin{itemize}
\raggedright
  \item \emph{Say it:} \emph{yogaḥ}
  \item \emph{Say it:} \emph{yogaḥ}, once more
  \item \emph{Say it:} say \emph{yogaḥ}
  \item \emph{Recall:} say the Sanskrit for a doctor, then the Sanskrit for exercise, then say \emph{yogaḥ} again
\end{itemize}

\begin{tcolorbox}[breakable,colback=teal!4,colframe=teal!35!black,title={Before you move on}]
What does \sk{योगः} mean? (\textquotedblleft{}Yoga\textquotedblright{}.) Say it once more.
\end{tcolorbox}
\section[\texorpdfstring{\sk{श्रमः} (śramaḥ)}{श्रमः (śramaḥ)}]{\sk{श्रमः} (śramaḥ) — toil, fatigue}

\label{lesson:SA-C211-sramah}



\begin{quote}
Before the new one: say the Sanskrit for exercise, then the Sanskrit for yoga.
\end{quote}

\subsection*{You'll want to know: \sk{श्रमः}}
\textbf{\sk{श्रमः}} — \emph{śramaḥ} — \textquotedblleft{}toil, fatigue\textquotedblright{}.

\begin{grammarlens}[title={What you've built}]
One more word for this chapter.
\end{grammarlens}

\subsection*{Your turn}
\begin{itemize}
\raggedright
  \item \emph{Say it:} \emph{śramaḥ}
  \item \emph{Say it:} \emph{śramaḥ}, once more
  \item \emph{Say it:} say \emph{śramaḥ}
  \item \emph{Recall:} say the Sanskrit for exercise, then the Sanskrit for yoga, then say \emph{śramaḥ} again
\end{itemize}

\begin{tcolorbox}[breakable,colback=teal!4,colframe=teal!35!black,title={Before you move on}]
What does \sk{श्रमः} mean? (\textquotedblleft{}Toil, fatigue\textquotedblright{}.) Say it once more.
\end{tcolorbox}
\section[\texorpdfstring{\sk{क्लान्तिः} (klāntiḥ)}{क्लान्तिः (klāntiḥ)}]{\sk{क्लान्तिः} (klāntiḥ) — tiredness}

\label{lesson:SA-C211-klantih}



\begin{quote}
Before the new one: say the Sanskrit for yoga, then the Sanskrit for toil.
\end{quote}

\subsection*{You'll want to know: \sk{क्लान्तिः}}
\textbf{\sk{क्लान्तिः}} — \emph{klāntiḥ} — \textquotedblleft{}tiredness\textquotedblright{}.

\begin{grammarlens}[title={What you've built}]
One more word for this chapter.
\end{grammarlens}

\subsection*{Your turn}
\begin{itemize}
\raggedright
  \item \emph{Say it:} \emph{klāntiḥ}
  \item \emph{Say it:} \emph{klāntiḥ}, once more
  \item \emph{Say it:} say \emph{klāntiḥ}
  \item \emph{Recall:} say the Sanskrit for yoga, then the Sanskrit for toil, then say \emph{klāntiḥ} again
\end{itemize}

\begin{tcolorbox}[breakable,colback=teal!4,colframe=teal!35!black,title={Before you move on}]
What does \sk{क्लान्तिः} mean? (\textquotedblleft{}Tiredness\textquotedblright{}.) Say it once more.
\end{tcolorbox}
\section[\texorpdfstring{\sk{स्नानम्} (snānam)}{स्नानम् (snānam)}]{\sk{स्नानम्} (snānam) — a bath}

\label{lesson:SA-C211-snanam}



\begin{quote}
Before the new one: say the Sanskrit for toil, then the Sanskrit for tiredness.
\end{quote}

\subsection*{You'll want to know: \sk{स्नानम्}}
\textbf{\sk{स्नानम्}} — \emph{snānam} — \textquotedblleft{}a bath\textquotedblright{}.

\begin{grammarlens}[title={What you've built}]
One more word for this chapter.
\end{grammarlens}

\subsection*{Your turn}
\begin{itemize}
\raggedright
  \item \emph{Say it:} \emph{snānam}
  \item \emph{Say it:} \emph{snānam}, once more
  \item \emph{Say it:} say \emph{snānam}
  \item \emph{Recall:} say the Sanskrit for toil, then the Sanskrit for tiredness, then say \emph{snānam} again
\end{itemize}

\begin{tcolorbox}[breakable,colback=teal!4,colframe=teal!35!black,title={Before you move on}]
What does \sk{स्नानम्} mean? (\textquotedblleft{}A bath\textquotedblright{}.) Say it once more.
\end{tcolorbox}
\section[(dialogue)]{Words from chapters 203-206 — a first pass}

\label{lesson:SA-R211-first-pass-words-from-chapters-203-206}



\begin{quote}
Say the words for tiredness and a bath.
\end{quote}

\subsection*{Your turn}
Say these aloud:

\begin{itemize}
\raggedright
  \item rejoices, salutes, begs, speaks, shines
  \item points out, digs, sows, is alive, heats
  \item forgives, blames, praises, accepts, blossoms
  \item sinks, bathes, frolics, shuts, opens
\end{itemize}

\begin{tcolorbox}[breakable,colback=teal!4,colframe=teal!35!black,title={Before you move on}]
Rejoices? (\textbf{\sk{मोदते}}, \emph{modate}.) Salutes? (\textbf{\sk{वन्दते}}, \emph{vandate}.) Begs? (\textbf{\sk{याचते}}, \emph{yācate}.) Speaks? (\textbf{\sk{भाषते}}, \emph{bhāṣate}.) Shines? (\textbf{\sk{शोभते}}, \emph{śobhate}.) Points out? (\textbf{\sk{सूचयति}}, \emph{sūcayati}.) Digs? (\textbf{\sk{खनति}}, \emph{khanati}.) Sows? (\textbf{\sk{वपति}}, \emph{vapati}.) Is alive? (\textbf{\sk{जीवति}}, \emph{jīvati}.) Heats? (\textbf{\sk{तपति}}, \emph{tapati}.) Forgives? (\textbf{\sk{क्षमते}}, \emph{kṣamate}.) Blames? (\textbf{\sk{निन्दति}}, \emph{nindati}.) Praises? (\textbf{\sk{प्रशंसति}}, \emph{praśaṁsati}.) Accepts? (\textbf{\sk{स्वीकरोति}}, \emph{svīkaroti}.) Blossoms? (\textbf{\sk{विकसति}}, \emph{vikasati}.) Sinks? (\textbf{\sk{मज्जति}}, \emph{majjati}.) Bathes? (\textbf{\sk{स्नाति}}, \emph{snāti}.) Frolics? (\textbf{\sk{खेलति}}, \emph{khelati}.) Shuts? (\textbf{\sk{निमीलयति}}, \emph{nimīlayati}.) Opens? (\textbf{\sk{उन्मीलयति}}, \emph{unmīlayati}.)
\end{tcolorbox}
\section[(dialogue)]{Words from chapters 207-211 — a second pass}

\label{lesson:SA-R211-second-pass-words-from-chapters-207-211}



\begin{quote}
Say the words for gambles and flows.
\end{quote}

\subsection*{Your turn}
Say these aloud:

\begin{itemize}
\raggedright
  \item gambles, flows, tastes, advises, honours
  \item a household, a relative, a companion, an enemy, a neighbour
  \item a poet, a singer, a dancer, a painter, a craftsman
  \item a preceptor, a head teacher, a pupil, a classmate, a doctor
  \item exercise, yoga, toil, tiredness, a bath
\end{itemize}

\begin{tcolorbox}[breakable,colback=teal!4,colframe=teal!35!black,title={Before you move on}]
Gambles? (\textbf{\sk{दीव्यति}}, \emph{dīvyati}.) Flows? (\textbf{\sk{प्रवहति}}, \emph{pravahati}.) Tastes? (\textbf{\sk{आस्वादयति}}, \emph{āsvādayati}.) Advises? (\textbf{\sk{उपदिशति}}, \emph{upadiśati}.) Honours? (\textbf{\sk{सम्मानयति}}, \emph{sammānayati}.) A household? (\textbf{\sk{परिवारः}}, \emph{parivāraḥ}.) A relative? (\textbf{\sk{बन्धुः}}, \emph{bandhuḥ}.) A companion? (\textbf{\sk{सखा}}, \emph{sakhā}.) An enemy? (\textbf{\sk{शत्रुः}}, \emph{śatruḥ}.) A neighbour? (\textbf{\sk{प्रतिवेशी}}, \emph{prativeśī}.) A poet? (\textbf{\sk{कविः}}, \emph{kaviḥ}.) A singer? (\textbf{\sk{गायकः}}, \emph{gāyakaḥ}.) A dancer? (\textbf{\sk{नर्तकः}}, \emph{nartakaḥ}.) A painter? (\textbf{\sk{चित्रकारः}}, \emph{citrakāraḥ}.) A craftsman? (\textbf{\sk{शिल्पी}}, \emph{śilpī}.) A preceptor? (\textbf{\sk{आचार्यः}}, \emph{ācāryaḥ}.) A head teacher? (\textbf{\sk{प्रधानाचार्यः}}, \emph{pradhānācāryaḥ}.) A pupil? (\textbf{\sk{विद्यार्थी}}, \emph{vidyārthī}.) A classmate? (\textbf{\sk{सहपाठी}}, \emph{sahapāṭhī}.) A doctor? (\textbf{\sk{चिकित्सकः}}, \emph{cikitsakaḥ}.) Exercise? (\textbf{\sk{व्यायामः}}, \emph{vyāyāmaḥ}.) Yoga? (\textbf{\sk{योगः}}, \emph{yogaḥ}.) Toil? (\textbf{\sk{श्रमः}}, \emph{śramaḥ}.) Tiredness? (\textbf{\sk{क्लान्तिः}}, \emph{klāntiḥ}.) A bath? (\textbf{\sk{स्नानम्}}, \emph{snānam}.)
\end{tcolorbox}
